\begin{table}[htbp]
\centering
\small
\caption{Distributional Incidence of the AI Grid Shock Across Household Income Quintiles}
\label{tab:table5_distribution}
\adjustbox{max width=\textwidth}{%
\begin{tabular}{lccccc}
\toprule
 & Quintile 1 & Quintile 2 & Quintile 3 & Quintile 4 & Quintile 5 \\
Dependent Variable: & (Bottom 20\%) & (20--40\%) & (40--60\%) & (60--80\%) & (Top 20\%) \\
\midrule
Post-Shock $\times$ Constrained LDA & 2.380*** & 1.060*** & 0.630*** & 0.360*** & 0.210*** \\
 & (0.065) & (0.038) & (0.024) & (0.016) & (0.011) \\
Pre-Shock Baseline Mean (\%) & 9.27 & 4.52 & 2.88 & 1.94 & 1.43 \\
Relative Burden Increase (\%) & 25.7\% & 23.5\% & 21.9\% & 18.6\% & 14.8\% \\
\midrule
County Fixed Effects & Yes & Yes & Yes & Yes & Yes \\
Quarter Fixed Effects & Yes & Yes & Yes & Yes & Yes \\
Observations & 3,600 & 3,600 & 3,600 & 3,600 & 3,600 \\
\bottomrule
\end{tabular}%
}
\begin{minipage}{\linewidth}
\vspace{1ex}
\footnotesize \textit{Notes:} Standard errors clustered at the county level in parentheses. * $p<0.10$, ** $p<0.05$, *** $p<0.01$. The dependent variable is household electricity expenditure as a percentage of gross income, calibrated from CEX and ACS microdata. Estimates document a strictly monotonic regressive gradient: low-income households incur a 2.38 percentage point increase in energy burden (+25.7\% relative expansion), compared to only 0.21 percentage points (+14.8\%) for top quintile earners, establishing a relative regressivity multiplier of $25.7\% / 14.8\% = 1.74\times$ and an absolute burden multiplier of $2.38 / 0.21 = 11.3\times$.
\end{minipage}
\end{table}
